\begin{theorem}[Lee--Yang 碰撞素数处幂次型分解的环面秩签名刚性]\label{thm:xi-terminal-zm-delta-s5-tame-collision-torus-rank-signature}
\leanverified{paper\_xi\_terminal\_zm\_delta\_s5\_tame\_collision\_torus\_rank\_signature}
沿用定理 \ref{thm:xi-terminal-zm-delta-s5-tame-collision-semistable-fiber} 的设定，并采用结论 \ref{con:xi-terminal-zm-delta-s5-jacobian-power-isogeny-refinement} 中的
\(
A_\varepsilon,E,B_3,B_5,B_6,B_7
\)
记号。
对在 \(p\) 处半稳定的阿贝尔簇 \(A/\QQ\)，记 \(r_p(A)\) 为其 N\'eron 模型单位连通分量的最大分裂环面秩。
则对每个 \(p\in P_2\)，六个同源因子 \(A_\varepsilon,E,B_3,B_5,B_6,B_7\) 均在 \(p\) 处半稳定，且其环面秩满足刚性签名
\[
r_p(E)=1,\qquad
r_p(B_3)=2,\qquad
r_p(B_5)=3,\qquad
r_p(B_6)=3,\qquad
r_p(B_7)=3,\qquad
r_p(A_\varepsilon)=1.
\]
因此其阿贝尔部分维数分别为
\[
\dim(E)-r_p(E)=0,\qquad
\dim(B_3)-r_p(B_3)=1,\qquad
\dim(B_5)-r_p(B_5)=2,\qquad
\dim(B_6)-r_p(B_6)=3,\qquad
\dim(B_7)-r_p(B_7)=4,\qquad
\dim(A_\varepsilon)-r_p(A_\varepsilon)=1.
\]
\end{theorem}
\begin{proof}
由推论 \ref{cor:xi-terminal-zm-delta-s5-tame-collision-neron-conductor}，\(J(X)\) 在每个 \(p\in P_2\) 处半稳定；
且其 \(S_5\)-等型半阿贝尔分裂由推论 \ref{cor:xi-terminal-zm-delta-s5-tame-collision-isotypic-neron-split} 控制。
另一方面，定理 \ref{thm:xi-terminal-zm-delta-s5-tame-collision-monodromy-induced-sgn} 给出单值化像的等型导子向量
\[
f_p(\varepsilon)=1,\quad
f_p((4,1))=4,\quad
f_p((3,2))=10,\quad
f_p((2,1,1,1))=12,\quad
f_p((2,2,1))=15,\quad
f_p((3,1,1))=18,
\]
并且在温和半稳定情形下，该等型维数恰等于对应等型半阿贝尔因子的环面秩。
故
\[
r_p(A_{(4,1)})=4,\quad
r_p(A_{(3,2)})=10,\quad
r_p(A_{(3,1,1)})=18,\quad
r_p(A_{(2,2,1)})=15,\quad
r_p(A_{(2,1,1,1)})=12,\quad
r_p(A_\varepsilon)=1.
\]
由结论 \ref{con:xi-terminal-zm-delta-s5-jacobian-power-isogeny-refinement}，
\[
A_{(4,1)}\sim_{\QQ}E^{\,4},\quad
A_{(3,2)}\sim_{\QQ}B_3^{\,5},\quad
A_{(3,1,1)}\sim_{\QQ}B_5^{\,6},\quad
A_{(2,2,1)}\sim_{\QQ}B_6^{\,5},\quad
A_{(2,1,1,1)}\sim_{\QQ}B_7^{\,4}.
\]
半稳定环面秩对同源不变且对直积可加，故
\(
r_p(E)=r_p(A_{(4,1)})/4=1
\)
以及其余各因子同理得到所列数值。
最后，半阿贝尔分解中维数可加，故阿贝尔部分维数为 \(\dim(A)-r_p(A)\)。
证毕。
\end{proof}

\begin{theorem}[几何--算术驯导子同一性：换位惯性下的指数匹配]\label{thm:xi-terminal-zm-delta-s5-tame-collision-geometry-arithmetic-conductor-identity}
沿用结论 \ref{con:xi-terminal-zm-delta-b5-discriminant-ramification-quadratic} 的 \(S_5\)-分裂域 \(L_B/\QQ\)。
取 \(p\in P_2\) 并取任一素理想 \(\mathfrak q\mid p\) 于 \(L_B\)。
则 \(p\nmid 120\) 且 \(B_5\bmod p\) 恰有一个二重根（结论 \ref{con:xi-terminal-zm-delta-b5-bad-reduction-primes-collision-fingerprint}），从而 \(\mathfrak q\) 处惯性群 \(I_{\mathfrak q}\) 为由一换位 \(\tau\in S_5\) 生成的温和循环群 \(C_2\)。
对任意不可约 \(S_5\)-表示 \(V_\lambda\)（特征标 \(\chi_\lambda\)），令 \(\rho_\lambda\) 为对应的 Artin 表示，则其驯 Artin 导子指数满足
\[
a_{\mathfrak q}(\rho_\lambda)=\dim V_\lambda-\dim\!\bigl(V_\lambda^{I_{\mathfrak q}}\bigr)
=\frac{\dim V_\lambda-\chi_\lambda(\tau)}{2}.
\]
并且在结论 \ref{con:xi-terminal-zm-delta-s5-jacobian-power-isogeny-refinement} 的各因子上，有完全可计算的等式
\[
r_p(A_\varepsilon)=a_{\mathfrak q}(\rho_{(1^5)}),\qquad
r_p(E)=a_{\mathfrak q}(\rho_{(4,1)}),\qquad
r_p(B_3)=a_{\mathfrak q}(\rho_{(3,2)}),\qquad
r_p(B_5)=a_{\mathfrak q}(\rho_{(3,1,1)}),\qquad
r_p(B_6)=a_{\mathfrak q}(\rho_{(2,2,1)}),\qquad
r_p(B_7)=a_{\mathfrak q}(\rho_{(2,1,1,1)}).
\]
\end{theorem}
\begin{proof}
温和情形下 \(I_{\mathfrak q}=\langle\tau\rangle\simeq C_2\)，故
\[
\dim V_\lambda^{I_{\mathfrak q}}
=\frac12\bigl(\chi_\lambda(1)+\chi_\lambda(\tau)\bigr)
=\frac12\bigl(\dim V_\lambda+\chi_\lambda(\tau)\bigr),
\]
从而
\(
a_{\mathfrak q}(\rho_\lambda)=\frac12(\dim V_\lambda-\chi_\lambda(\tau))
\)。
另一方面，由定理 \ref{thm:xi-terminal-zm-delta-s5-tame-collision-monodromy-induced-sgn}，
在几何侧的单值化像满足 \(\mathrm{Im}(N)\cong \Ind_{\langle\tau\rangle}^{S_5}(\mathrm{sgn})\)，且对任意不可约 \(V_\lambda\) 有
\[
m_\lambda
=\dim\Hom_{\langle\tau\rangle}(V_\lambda,\mathrm{sgn})
=\frac12\bigl(\dim V_\lambda-\chi_\lambda(\tau)\bigr).
\]
由定理 \ref{thm:xi-terminal-zm-delta-s5-tame-collision-torus-rank-signature} 的证明可知：
温和半稳定退化下，\(m_\lambda\) 恰等于结论 \ref{con:xi-terminal-zm-delta-s5-jacobian-power-isogeny-refinement} 中对应因子的环面秩 \(r_p\)。
因此 \(r_p(\cdot)=m_\lambda=a_{\mathfrak q}(\rho_\lambda)\) 于各因子成立。
证毕。
\end{proof}

\endinput
